\scanpage{17}

\noindent\hfill\unsure{Villecun|Villecui} le 23.6.1974\\

Chers Deligne, Giraud, Verdier,\\

Vous savez peut-être qu'une mathématicienne vietnamienne, Hoang Xuan Sinh, \inserted{a} travaillé depuis mon séjour au Vietnam en 1967 à un sujet de recherche que je lui avais proposé, savoir une théorie de structure des (Gr)-catégories, et la construction d\deleted{\unsure{b|h|n}}'une (Gr)-catégorie (ou plutôt catégorie de Picard) enveloppante pour une catégorie avec produit $\unsure{\otimes|\&}$ associatif et commutatif. J'ai reçu récemment une copie du manuscript en forme (dont la confection s'est faite au prix de mille difficultés - il est douteux qu'il soit possible d'avoir d'autres copies dans un avenir proche, à moins d'en faire d'après cet exemplaire-là), et j'ai envoyé un rapport officiel au sujet de ce travail à M. Ta Quang B$\overset{\handwritten{\unsure{\,)}}}{\mathrm{u}}\handwritten{\breve}{\mathrm{u}}$, ministre de l'enseignement supérieur et technique de la RDV, rapport dont je vous envoie une copie ci-joint. Mme Sinh, avec mon accord, a demandé à présenter ce travail comme thèse de doctorat, et à faire sa soutenance en France. \deleted{\unsure{I|T}}A l'heure où j'écris, je n'ai pas connaissance d'un accord des autorités vietnamiennes pour l'un ni pour l'autre; il y a eu un certain nombre de difficultés soulevées, qu'il est intuile de détailler ici, et qui j'espère vont s'aplanir. J'ai appris par ailleurs que Mme Hoang \unsure{\deleted{H}X|\deleted{X}H}uan Sinh figurerait sur la liste officielle de la délégation vietnamienne au congrès international de mathématiciens à Vancouver du 20 août 1974. Si c'est bien ainsi, il serait très souhaitable qu'elle puisse profiter de cette occasion pour faire sa soutenance à Orsay au retour de Vancouver, fin Septembre ou début Octobre, et il me semble plausible que les autorités vietnamiennes y donneront leur accord. Aussi il serait assez urgent, vu les lenteurs administratives, de constituer un jury. Par la présente, je vous demande si vous seriez disposé en principe à faire partie (avec moi) de ce jury de thèse - ce qui suppose évidemment que vous vous trouvez dans la région parisienne au moment prévu pour la soutenance. A celui qui m'en ferait la demande, je peux envoyer le texte de la thèse complète, ou un résumé assez détaillé rédigé par Mme Hoang Xuan Sinh , ou enfin un résumé moins détaillé de 13 pages manuscrites (rédigées en vue d'une publication aux CR, mais beaucoup trop long à cet effet). Ce serait gentil de me répondre rapidement, et en cas d'impossibilité, de voir autour de vous s'il y a des gens ``dans le bain'' d\deleted{\unsure{e|a}}\handwritten{u} \deleted{\unsure{ce}}\handwritten{genre} \handwritten{\inserted{de choses}} que fait Mme Sinh, et qui seraient intéressés à faire partie du jury.
% archive sheet number: 16
